\documentclass[11pt,a4paper]{article}
\usepackage[margin=27mm]{geometry}
\usepackage{amsmath,amssymb,amsthm,booktabs,array}
\usepackage[T1]{fontenc}
\usepackage{lmodern}
\usepackage[colorlinks=true,linkcolor=blue,citecolor=blue,urlcolor=blue]{hyperref}
\newcommand{\qbin}[2]{\genfrac{[}{]}{0pt}{}{#1}{#2}_{q}}
\newcommand{\coeff}[2]{[q^{#1}]#2}
\newtheorem{theorem}{Theorem}
\newtheorem{corollary}[theorem]{Corollary}
\title{Proofs of Okada's Conjectures 9.63 and 9.64\\with constant-base-shift extensions\\[3pt]
\large arXiv:2003.02112v2, Section 9.8.3}
\author{Haobo Ma$^{1}$ \qquad Wenlin Zhang$^{2,1}$\\[4pt]
\small $^{1}$ChronoAI Pte Ltd, Singapore\\
\small $^{2}$National University of Singapore}
\date{November 3, 2026}
\begin{document}
\maketitle
\begin{abstract}
We prove Conjectures 9.63 and 9.64 in
Section 9.8.3 of Okada's arXiv:2003.02112v2: every almost strictly pre-fitting, unit-weight,
unit-base-shift pre-merged Gaussian determinant of width at least three
and degree shift zero fails to be almost palindromic, and the proposed
width-two proportions are strictly increasing from index three onward.
The second nonzero coefficient exceeds the penultimate nonzero coefficient
by one. At width two, almost palindromicity holds exactly when one of the
two tuple gaps equals one. This classification gives the exact proportion
$O(\lambda)=(\lambda-2)(\lambda-3)/(\lambda(\lambda+1))$ for $\lambda\ge2$.
These conclusions extend to every constant positive integral base shift
at degree shift zero. The proposed base-shift invariance at nonzero degree
shift can fail at width two, as in Okada's Example 9.65; an infinite
family extends that example. We prove invariant slices of Conjecture 9.66,
whose width is at least three.
\end{abstract}

\section{The conjectures and the coefficient comparison}

Okada's Conjecture 9.63 concerns the symmetry remaining after zero
coefficients are removed from a pre-merged determinant. We prove it by
comparing four coefficients, and obtain a width-two classification that
also proves the neighboring Conjecture 9.64. Throughout, conjecture and
definition numbers refer to the January 2024 version~\cite{okada}.
We also extend Example 9.65 to an infinite family and prove
base-shift invariance on several slices of Conjecture 9.66. These numbered
statements belong to v2; the author's November 2, 2026 manuscript~\cite{okada-current}
reorganizes the conjectures and does not repeat these three statements.

For $0\le k\le n$, write
\[
 Q_n(q)=\prod_{j=1}^{n}(1-q^j),\qquad
 \qbin{n}{k}=\frac{Q_n}{Q_kQ_{n-k}}\in\mathbb Z[q],\qquad Q_0=1.
\]
These are the standard Gaussian polynomials; see, for example,
Stanley~\cite[p.~501, equations (2)--(3)]{stanley}.
For natural tuples satisfying
\begin{equation}\label{fitting}
 0\le a_1<\cdots<a_\mu\le b_1<\cdots<b_\mu,
\end{equation}
put
\begin{equation}\label{products}
 A(q)=\prod_{i=1}^{\mu}\qbin{b_i}{a_i},\qquad
 B(q)=\prod_{i=1}^{\mu}\qbin{b_i}{a_{\mu+1-i}},\qquad P(q)=A(q)-B(q).
\end{equation}
Condition~\eqref{fitting} is exactly ``almost strictly pre-fitting'' in
\cite[Definition 9.48(2)]{okada}. At weights and base shifts all equal to
one, the base-shift factors in \cite[Definition 9.4]{okada} are one.
Its degree-shift-zero determinant is therefore precisely~\eqref{products}.
Following \cite[Definition 9.52]{okada}, a nonzero polynomial is
\emph{almost palindromic} when its nonzero coefficients, in increasing
degree, form a palindrome.

\begin{theorem}\label{main}
Assume~\eqref{fitting} and $\mu\ge3$. Define
\[
 t=b_1-a_\mu,\qquad
 c=\begin{cases}\mu,&a_1>0,\\ \mu-1,&a_1=0,\end{cases}
 \qquad d=\sum_{i=1}^{\mu}a_i(b_i-a_i).
\]
Then $\deg P=d$, $d\ge t+4$, and
\begin{equation}\label{endpoints}
 \coeff{j}{P}=0\ (j<t+1),\quad
 \coeff{t+1}{P}=1,\quad \coeff{t+2}{P}=c+1,\quad
 \coeff{d-1}{P}=c,\quad \coeff{d}{P}=1.
\end{equation}
In particular, $P$ is not almost palindromic, proving Conjecture 9.63.
The same conclusions hold at width two if
$a_2-a_1\ge2$ and $b_2-b_1\ge2$.
\end{theorem}

\begin{proof}
We prove both regimes together. For $n\ge1$, Gaussian polynomials satisfy
\[
 \qbin{n}{k}=\qbin{n-1}{k}+q^{n-k}\qbin{n-1}{k-1},
 \qquad \qbin{n}{0}=\qbin{n}{n}=1,
\]
with terms outside $0\le k\le n$ equal to zero. Thus they are integer
polynomials. Their factorial formula gives constant coefficient one and
\[
 \left.\qbin{n}{k}\right|_{q\mapsto q^{-1}}=q^{-k(n-k)}\qbin{n}{k}.
\]
Here the left side means evaluation of the same Gaussian polynomial at
$q^{-1}$. Consequently it is monic of degree $k(n-k)$.
If $0<k<n$, then $Q_n\equiv1-q\pmod{q^2}$ and
$Q_kQ_{n-k}\equiv1-2q\pmod{q^2}$, so its linear coefficient is one.
At $k=0$ or $k=n$ that coefficient is zero.

First compare the low degrees. For $\mu\ge3$, strict increase implies
\begin{align*}
 b_i-a_i&=t+(b_i-b_1)+(a_\mu-a_i)\ge t+\mu-1\ge t+2,\\
 b_1-a_\mu&=t,\\
 b_i-a_{\mu+1-i}&\ge t+2(i-1)\ge t+2\qquad(i\ge2).
\end{align*}
At width two with gaps $x=a_2-a_1\ge2$ and $y=b_2-b_1\ge2$,
the direct differences are $t+x,t+y$ and the reversed ones are $t,t+x+y$,
so the same bounds apply.

Every $Q_n$ is a unit in $\mathbb Z[[q]]$. Cancellation of the common
Gaussian factorials gives
\[
 A\prod_i Q_{b_i-a_i}=B\prod_i Q_{b_i-a_{\mu+1-i}}.
\]
In $\mathbb Z[[q]]/(q^{t+3})$, all $Q_n$ with $n\ge t+2$ equal
$Q_{t+2}$. The bounds just established therefore imply
\[
 B\equiv A\frac{Q_{t+2}}{Q_t}
   =A(1-q^{t+1})(1-q^{t+2})\pmod{q^{t+3}}.
\]
Since $2t+3\ge t+3$, subtraction yields
\begin{equation}\label{low}
 P\equiv A(q^{t+1}+q^{t+2})\pmod{q^{t+3}}.
\end{equation}
All direct differences are positive. Exactly $c$ factors of $A$ have
positive lower index, so $A(0)=1$ and $\coeff{1}{A}=c$.
Equation~\eqref{low} proves the three low-degree assertions
in~\eqref{endpoints}.

For the high degrees, set
$e=\sum_i a_{\mu+1-i}(b_i-a_{\mu+1-i})=\deg B$.
Pairing reversed indices gives the exact identity
\begin{equation}\label{gap}
 d-e=\sum_{i=1}^{\lfloor\mu/2\rfloor}
       (a_{\mu+1-i}-a_i)(b_{\mu+1-i}-b_i).
\end{equation}
Every summand is positive. For $\mu\ge3$, the outermost one is at least
$(\mu-1)^2\ge4$; in the width-two regime it equals $xy\ge4$.
Hence $\deg B\le d-4$. Since $A$ is monic and reciprocal of degree $d$,
its last two coefficients are $c,1$, and $B$ changes neither.
This proves $\deg P=d$ and the last two assertions of~\eqref{endpoints}.

Finally, in both regimes $a_\mu\ge2$ and
$b_\mu-a_\mu\ge t+2$. Thus
$d\ge a_\mu(b_\mu-a_\mu)\ge2(t+2)\ge t+4$.
The second and penultimate nonzero coefficients occupy the distinct
degrees $t+2$ and $d-1$. They are $c+1$ and $c$, respectively, and
$c\ge1$. Removing any zeros between them cannot make the list palindromic.
\end{proof}

For example, $a=(0,1,2)$ and $b=(2,3,4)$ give
$P=q+3q^2+4q^3+4q^4+2q^5+q^6$; the second and penultimate coefficients
are three and two.

\section{Width two and the exact counting formula}

The width-two sentence immediately following Conjecture 9.63 states the
nonpalindromicity conclusion for two gaps at least two.
Theorem~\ref{main} proves the stronger conclusion of failure of almost
palindromicity. The remaining cases have a simple exact description.

\begin{theorem}\label{classification}
For $0\le a_1<a_2\le b_1<b_2$, the polynomial~\eqref{products} at width
two is almost palindromic if and only if
\[
 \min(a_2-a_1,b_2-b_1)=1.
\]
\end{theorem}
\begin{proof}
If both gaps are at least two, apply Theorem~\ref{main}.
Otherwise put $x=a_2-a_1$, $y=b_2-b_1$, $r=b_1-a_2+1$, and
$h=\max(x,y)$. Direct cancellation gives
\[
 \frac BA=
 \frac{Q_{b_1-a_1}Q_{b_2-a_2}}{Q_{b_1-a_2}Q_{b_2-a_1}}
 =\frac{1-q^r}{1-q^{r+h}}.
\]
For $x=1$, cancel to obtain
$(1-q^{b_1-a_2+1})/(1-q^{b_2-a_1})$; for $y=1$, the same
expression follows by canceling the other two adjacent factorials.
Both denominators have exponent $r+h$. Thus, in $\mathbb Q(q)$,
\begin{equation}\label{onegap}
 P=q^r A\frac{1-q^h}{1-q^{r+h}}.
\end{equation}
The right side has order $r$ and coefficient one there. Moreover,
$d-e=xy=h>0$, so $\deg P=d$.
The rational factor $H(q)=(1-q^h)/(1-q^{r+h})$ satisfies
$H(q^{-1})=q^rH(q)$. Combining this with reciprocity of $A$ in
\eqref{onegap} gives
\[
 q^{d+r}P(q^{-1})=P(q).
\]
The coefficients on the interval $[r,d]$ are therefore symmetric,
including all zero coefficients. Deleting zeros preserves that symmetry.
\end{proof}

\begin{corollary}\label{counting}
Let $M(\lambda)$ and $N(\lambda)$ be the sets in Okada's Conjecture 9.64:
$M(\lambda)$ consists of all $0\le a_1<a_2\le b_1<b_2=\lambda$, and
$N(\lambda)$ consists of those tuples whose determinant is not almost
palindromic. For $\lambda\ge2$,
\[
 |M(\lambda)|=\binom{\lambda+1}{3},\quad
 |N(\lambda)|=\binom{\lambda-1}{3},\quad
 |M(\lambda)\setminus N(\lambda)|=(\lambda-1)^2,
\]
where a binomial coefficient with upper index less than three is zero.
Consequently
\begin{equation}\label{ratio}
 O(\lambda)=\frac{|N(\lambda)|}{|M(\lambda)|}
           =\frac{(\lambda-2)(\lambda-3)}{\lambda(\lambda+1)}.
\end{equation}
In particular, $O(3)=0$, $O(\lambda+1)>O(\lambda)$ for $\lambda\ge3$,
and $O(\lambda)\longrightarrow1$.
\end{corollary}
\begin{proof}
The bijection
\[
 (a_1,a_2,b_1,\lambda)=(u,u+x,u+x+t,u+x+t+y)
\]
identifies $M(\lambda)$ with $u,t\ge0$, $x,y\ge1$, and
$u+t+x+y=\lambda$. Subtracting one from $x$ and $y$ leaves four
nonnegative integers of sum $\lambda-2$, giving $\binom{\lambda+1}{3}$
by stars and bars. Theorem~\ref{classification} identifies $N(\lambda)$
with $x,y\ge2$. Subtracting two instead gives $\binom{\lambda-1}{3}$
for $\lambda\ge4$, and no tuples for $\lambda=2,3$.
Subtracting the two counts gives $(\lambda-1)^2$, and division gives
\eqref{ratio}. Finally,
\[
 O(\lambda+1)-O(\lambda)
 =\frac{6(\lambda-2)}{\lambda(\lambda+1)(\lambda+2)}>0\quad(\lambda\ge3).
\]
The limit follows immediately from~\eqref{ratio}.
The ratio is undefined at $\lambda=0,1$, where $M(\lambda)$ is empty;
at $\lambda=2$ it equals zero.
\end{proof}

\begin{corollary}\label{allwidths}
At fixed $b_\mu=\lambda$ and width $\mu\ge3$, the number of tuples
in~\eqref{fitting} is $\binom{\lambda+1}{2\mu-1}$, and every one gives
a determinant that is not almost palindromic. At width two the exceptional
proportion is $1-O(\lambda)=6(\lambda-1)/(\lambda(\lambda+1))$,
asymptotic to $6/\lambda$.
\end{corollary}
\begin{proof}
If $a_\mu<b_1$, choose the other $2\mu-1$ distinct entries from
$\{0,\ldots,\lambda-1\}$. If $a_\mu=b_1$, choose $2\mu-2$ entries
from that set. Pascal's identity gives the stated sum, and
Theorem~\ref{main} applies. The last formula follows from~\eqref{ratio}.
\end{proof}

\section{Constant base shifts and neighboring statements}

The source's base-shift function is
\[
 b_{n,\rho}(q)=\frac{Q_n(q^\rho)}{Q_n(q)}
 =\prod_{j=1}^{n}(1+q^j+\cdots+q^{(\rho-1)j}),\qquad \rho\ge1.
\]
In Definition 9.4 at unit weights and constant base shift $\rho$, set
\begin{align*}
 A_\rho&=\prod_i\left(\qbin{b_i}{a_i}\bigg|_{q\mapsto q^\rho}
                b_{b_i-a_i,\rho}(q)\right),\\
 B_\rho&=\prod_i\left(\qbin{b_i}{a_{\mu+1-i}}\bigg|_{q\mapsto q^\rho}
                b_{b_i-a_{\mu+1-i},\rho}(q)\right),\qquad
 P_{\rho,e}=A_\rho-q^eB_\rho.
\end{align*}
Thus $A_1=A$, $B_1=B$, and $P_{1,0}=P$; changing $\rho$ is not just
substituting $q^\rho$ into $P$.

\begin{corollary}\label{flatbase}
The non-almost-palindromicity conclusion of Theorem~\ref{main}, the
width-two classification of Theorem~\ref{classification}, and all counts
above hold for $P_{\rho,0}$ for every integer $\rho\ge1$.
In the regimes of Theorem~\ref{main}, its first two nonzero coefficients
are $1,c_\rho+1$ and its last two are $c_\rho,1$, where
$c_1=c$ and $c_\rho=\mu$ for $\rho\ge2$.
\end{corollary}
\begin{proof}
Cancelling the factorials in each factor gives
\[
 \qbin{n}{k}\bigg|_{q\mapsto q^\rho} b_{n-k,\rho}(q)
 =\frac{Q_n(q^\rho)}{Q_k(q^\rho)Q_{n-k}(q)}.
\]
Consequently $B_\rho/A_\rho=B/A$ as rational functions.
Each factor of $A_\rho$ is a monic reciprocal polynomial, with degree
\[
 D_\rho(n,k)=\rho k(n-k)+\frac{\rho-1}{2}(n-k)(n-k+1).
\]
Write $d_\rho=\sum_iD_\rho(b_i,a_i)$.
The difference of the direct and reversed degrees is the same positive
integer $G=\deg A-\deg B$ in~\eqref{gap}: the Gaussian contributions give
$\rho G$, while the base-shift contributions give $-(\rho-1)G$.
If $\rho\ge2$ and $n-k>0$, the Gaussian factor has no linear term
and $b_{n-k,\rho}$ has linear coefficient one. Thus
$[q]A_\rho=c_\rho$. The congruence~\eqref{low} and the high-degree
comparison apply with $A_\rho$ and $d_\rho$; also $d_\rho\ge d$.
For width two with a unit gap, the rational factor
in~\eqref{onegap} is unchanged. Reciprocity now yields
$q^{d_\rho+r}P_{\rho,0}(q^{-1})=P_{\rho,0}(q)$.
This proves both cases of the classification and hence the counts.
\end{proof}

\begin{corollary}\label{weighted}
Assume~\eqref{fitting} with $\mu\ge2$. For every constant base shift $\rho\ge1$ and
every constant integral weight $w\ge2$, the degree-shift-zero determinant
$A_\rho^w-B_\rho^w$ is not almost palindromic.
\end{corollary}
\begin{proof}
Put $t=b_1-a_\mu$. The direct differences are at least $t+1$,
and all reversed differences except the first are at least $t+2$.
Hence $B_\rho/A_\rho\equiv1-q^{t+1}\pmod{q^{t+2}}$.
The first nonzero coefficient of $A_\rho^w-B_\rho^w$ is therefore $w$,
whereas its last is one because its direct degree exceeds its reversed
degree by $wG>0$.
\end{proof}

Conjecture 9.66 proposes that almost palindromicity is unchanged when
the constant base shift increases from $\rho$ to $\rho+1$, for
\emph{widths at least three} and degree shifts $0\le e\le G$.
At width two the assertion is false, as Example 9.65 already shows for
$a=(0,1)$, $b=(1,4)$ and degree shift one. The next corollary extends
that example; it is not a counterexample to Conjecture 9.66.

\begin{corollary}\label{counterexample}
For every $n\ge4$, let $a=(0,1)$, $b=(1,n)$ and $e=1$.
Then $P_{1,1}$ is almost palindromic, whereas $P_{\rho,1}$ is not
almost palindromic for every $\rho\ge2$.
\end{corollary}
\begin{proof}
Here $G=n-1$, so $e=1$ is admissible, and
\[
 P_{1,1}=1+q^2+q^3+\cdots+q^{n-1}.
\]
Its nonzero coefficients are all one. For $\rho\ge2$, the linear
coefficient of $A_\rho$ is two and $B_\rho(0)=1$, so
$[q]P_{\rho,1}=1$ and $P_{\rho,1}(0)=1$.
The degree of $qB_\rho$ is $d_\rho-(n-2)\le d_\rho-2$.
Thus the last two coefficients are $2,1$, by reciprocity of $A_\rho$.
The second and penultimate nonzero coefficients differ.
\end{proof}

For the smallest member, $n=4$, direct multiplication gives
\begin{align*}
 P_{1,1}&=1+q^2+q^3,\\
 P_{2,1}&=1+q+2q^2+4q^3+4q^4+5q^5+6q^6+6q^7\\
         &\hspace{12mm}+5q^8+5q^9+4q^{10}+2q^{11}+2q^{12}+q^{13}.
\end{align*}

These are precisely the two polynomials in Example 9.65.
The first is \emph{almost} palindromic: its coefficient list including
the interior zero is $(1,0,1,1)$, which is not palindromic, whereas
its nonzero list is $(1,1,1)$. Thus the word ``palindromic'' there
requires this qualification.

\begin{corollary}\label{shifted}
Assume~\eqref{fitting} with $\mu\ge2$. For any $\rho\ge1$, put
$G=\deg A_\rho-\deg B_\rho$.
Then
\[
 q^{d_\rho}P_{\rho,G}(q^{-1})=P_{\rho,0}(q).
\]
Hence the zero-shift classification also holds at $e=G$.
If $\mu\ge3$, the determinants at $e=1$ and $e=G-1$ are not almost
palindromic for any $\rho\ge1$.
If $G$ is even, $P_{\rho,G/2}$ is reciprocal, and therefore almost
palindromic. Multiplication by $q^s$ or substitution $q\mapsto q^k$,
for integers $s\ge0$, $k\ge1$, preserves almost palindromicity.
\end{corollary}
\begin{proof}
For $0\le e\le G$, reciprocity of $A_\rho$ and $B_\rho$ gives
$q^{d_\rho}P_{\rho,e}(q^{-1})=P_{\rho,G-e}(q)$.
For $e=G$ this reverses the full coefficient list, so it also reverses
the nonzero list. For $e=G/2$ it gives symmetry directly; the constant
coefficient is one. When $\mu\ge3$, $G\ge4$ and $c_\rho\ge2$.
At $e=1$ the constant and linear coefficients are $1,c_\rho-1$;
the last two are $c_\rho,1$, since $\deg(qB_\rho)\le d_\rho-3$.
Thus the second and penultimate nonzero coefficients differ.
Reciprocity gives the assertion at $e=G-1$.
The final two transformations preserve the ordered
nonzero coefficient list. An overall monomial shift is distinct from
the degree-shift parameter $e$, which affects the reversed product alone.
\end{proof}

The sentence before Conjecture 9.64 also anticipates infinitely many
almost palindromic unimodal polynomials. An explicit family is
$a=(0,1)$, $b=(m,n)$ with $1\le m<n$:
\[
 P_{1,0}=q^m(1+q+\cdots+q^{n-m-1}).
\]
This follows by subtracting $\qbin{m}{1}$ from $\qbin{n}{1}$.
Its nonzero coefficient sequence is constant, so it is palindromic and
unimodal. No general unimodality assertion is needed for this family.

\section{Coverage and mathematical verification}

The following table records the scope relative to the posing source~\cite{okada}.
Theorems~\ref{main}--\ref{classification} concern unit weights, unit base
shifts and degree shift zero; the corollaries state their extended ranges.

\begin{center}
\small
\begin{tabular}{@{}p{0.32\textwidth}p{0.60\textwidth}@{}}
\toprule
Source item & Coverage\\
\midrule
Conjecture 9.63 & Proved by Theorem~\ref{main}, for every width at least three.\\[3pt]
Sentence following 9.63 & Proved in the width-two part of Theorem~\ref{main};
Theorem~\ref{classification} gives both directions.\\[3pt]
Conjecture 9.64 & Proved by Corollary~\ref{counting}, including the exact counts and limit.\\[3pt]
Example 9.65 & Prior proof cited and wording corrected: the first
polynomial is almost palindromic, not palindromic including its interior zero.
Corollary~\ref{counterexample} extends it to all $n\ge4$, $\rho\ge2$.\\[3pt]
Conjecture 9.66 & Proved on slices $e=0,1,G-1,G$ and, for even $G$,
$e=G/2$, by Corollaries~\ref{flatbase} and~\ref{shifted}.
Other degree shifts are not covered; its hypothesis is $\mu\ge3$.\\[3pt]
Sentence before 9.64 & Proved for the explicit infinite unimodal family above.\\[3pt]
Conjectures 9.49, 9.53, 9.55, 9.57, 9.60, 9.62 & Not covered: density,
log-concavity, and unimodality require information beyond these endpoint coefficients.\\[3pt]
2026 manuscript, Conjectures 11.2, 11.4, 11.6 & Not covered: these concern
unimodality and log-concavity, rather than the palindromicity statements of v2.\\
\bottomrule
\end{tabular}
\end{center}

The pinned Lean source, dated November 2, 2026,
\href{https://github.com/the-omega-institute/trureturing/blob/0b3d3a2945e69db04f778eb6fea96ee6a8030a45/D5/S3/Combinatorics/GaussianQuotient/OkadaPreMergedNonPalindromic.lean}{\texttt{OkadaPreMergedNonPalindromic.lean}}
contains the declaration \texttt{result : claim}, which proves the
width-at-least-three conclusion of Theorem~\ref{main} for all natural
tuples~\eqref{fitting}. Its definitions \texttt{preMerged} and
\texttt{IsAlmostPalindromic} are respectively~\eqref{products} and the
palindrome condition on the ordered nonzero coefficients.
The lemmas \texttt{low\_end} and \texttt{high\_end} prove the
endpoint formulas~\eqref{endpoints} at width at least three and $\rho=1$;
the paired degree-gap identity~\eqref{gap}, width-two classification,
counting formulas, constant-base-shift and constant-weight results,
the family extending Example 9.65 and the stated slices of Conjecture 9.66
have the written proofs above.

The accompanying \texttt{verify.py} performs exact integer-polynomial
checks: 66 Gaussian polynomials with $0\le k\le n\le10$ are compared
with subset enumeration; all 1,486 fitting pairs with $3\le\mu\le6$
and $b_\mu\le10$ satisfy the endpoint formulas; and all 4,845 width-two
tuples with $b_2\le18$ satisfy Theorem~\ref{classification} and the counts
in Corollary~\ref{counting}. Additional exact checks cover constant base
shifts, weights, degree-shift reciprocity and the family extending Example 9.65;
their explicit finite ranges are recorded in \texttt{checked-results.json}.
These finite calculations supplement the all-parameter proofs.

\begin{thebibliography}{3}
\bibitem{okada}
So Okada, \emph{Merged-log-concavity of rational functions, almost strictly
unimodal sequences, and phase transitions of ideal boson-fermion gases},
\href{https://arxiv.org/abs/2003.02112v2}{arXiv:2003.02112v2},
January 2024, Definitions 9.4, 9.48, 9.52 and Section 9.8.3,
Conjectures 9.63, 9.64, Example 9.65 and Conjecture 9.66, pp.~96--97.

\bibitem{okada-current}
So Okada, \emph{Merged-log-concavity of rational functions, almost strictly
unimodal sequences, and phase transitions of ideal boson and fermion gases},
\href{https://so-okada.github.io/research_files/phases.pdf}{author-hosted manuscript},
November 2, 2026, Section 11, pp.~100--101.

\bibitem{stanley}
Richard P. Stanley, \emph{Log-concave and unimodal sequences in algebra,
combinatorics, and geometry}, in \emph{Graph Theory and Its Applications:
East and West}, Annals of the New York Academy of Sciences \textbf{576}
(1989), 500--535.
\href{https://doi.org/10.1111/j.1749-6632.1989.tb16434.x}{doi:10.1111/j.1749-6632.1989.tb16434.x}.
\end{thebibliography}

\medskip
\noindent\small Licensed under
\href{https://creativecommons.org/licenses/by/4.0/}{Creative Commons Attribution 4.0 International (CC BY 4.0)}.
\end{document}
